\begin{tomtat}
\begin{itemize}
  \item Khi lời giải thích do chính mô hình in ra cùng đáp án, vế trái của
định nghĩa~\ref{def:ch01-do-trung-thuc} đổi chủ, và bốn thứ đổi theo: không còn
tập can thiệp nào để công bố, lời giải thích và đáp án ra đời từ cùng một quá
trình nên vừa mô tả vừa gây ra, lời giải thích là văn bản nên khẳng định được
những điều không có thật, và đối tượng đã đổi một lần rồi. Bài neo xếp chuỗi
few-shot tuyến tính nguyên bản vào cùng lớp với bản đồ chú ý và xếp chuỗi của mô
hình suy luận ra ngoài lớp ấy, nên hỏi một chuỗi có trung thực không thì phải hỏi
thêm chuỗi của loại mô hình nào.
  \item Bài đưa chuỗi suy luận vào ngành đã ba lần từ chối tuyên bố nó giải thích
được mô hình: phần trong ngoặc của tính chất thứ hai nói việc mô tả đầy đủ các
tính toán nâng đỡ một đáp án vẫn là câu hỏi mở, bài chọn động từ
\enquote{bắt chước} thay vì gọi chuỗi là lời giải thích, và phụ lục C.2 gọi khả
năng diễn giải là hệ quả phụ chứ không phải mục tiêu. Trong 43 trang, chữ
\enquote{faithfulness} xuất hiện đúng hai lần, một lần trong một nhan đề được
trích và một lần trong danh sách các năng lực trồi lên theo quy mô, còn chữ
\enquote{unfaithful} không lần nào.
  \item Khảo sát về mô hình ngôn ngữ cho XAI xếp chuỗi suy luận vào nhóm
\tn{diễn giải nội tại}{intrinsic interpretability},
tức nhóm tự diễn giải được nên không phải đối chiếu, và dẫn cho
nước đi ấy một khảo sát khác chứ không phải bài mở đường, dù bài mở đường có trong
danh mục tài liệu của nó. Ô nguồn của dòng bảng ấy dẫn bài đo độ trung thực của
chuỗi suy luận, và bài được dẫn kết luận rằng mô hình càng lớn và càng mạnh thì
suy luận sinh ra càng kém trung thực, và rằng chuỗi \emph{có thể} trung thực nếu
chọn cẩn thận mô hình và nhiệm vụ. Số hiệu ấy xuất hiện đúng hai lần trong cả bài
khảo sát, một lần trong ô bảng và một lần trong danh mục tài liệu.
  \item Cùng khảo sát ấy đặt độ trung thực vào cột định lượng và không đưa ra
công thức, ngưỡng, thủ tục hay điểm số nào để tính nó trong cả 18 trang; ô ví dụ
của dòng ấy mô tả một hệ thống chứ không phải một con số, chữ \enquote{unfaithful}
không xuất hiện lần nào, và năm khó khăn mà mục hạn chế liệt kê không có khó khăn
nào là việc không ai kiểm được điều ấy. Bài của Jacovi và Goldberg thì được xếp ở
dòng tính hợp lý trong bảng, trong khi phần thân dẫn nó đúng một lần và dẫn đúng.
Mười hai con số duy nhất của bài, ở bảng~\ref{tab:ch17-saliency}, đo một quan hệ
khác, gọi tên một mô hình khác với hình của chính nó, và không đơn điệu theo quy
mô ở đúng chỗ hai mô hình bị xếp chung nhóm.
  \item Khảo sát 2026 về suy luận nêu chỗ nghi ba lần, vẽ hình 27 đặt độ trung
thực dưới mặt nước, và gọi tên ba họ dụng cụ, cả ba đều nằm trong bốn chỉ số mà
bảng~\ref{tab:ch09-auroc} đã đem ra kiểm. Nhưng câu định nghĩa trong hình ấy hỏi
chuỗi suy luận có gây ra đáp án hay không, tức là hỏi mức quan trọng chứ không
hỏi độ trung thực, và đó là phép lẫn mà chẩn đoán trung tâm của bài neo nói tới,
lần này nằm ngay trong một định nghĩa được in ra.
  \item Chiều của phép chuyển ngược với chiều mà chữ \enquote{chuyển} của
mục~\ref{sec:ch01-luan-de-cua-sach} gợi ra, và bằng chứng nằm trong sách:
mục~\ref{sec:ch09-doi-tuong-khac-nhau} đã ghi rằng ở phía
chuỗi suy luận câu hỏi đã được hỏi bằng thực nghiệm và câu trả lời là không, còn
ở phía attribution đặc trưng thì chưa ai hỏi. Phép đối chiếu không đi từ Phần~IV
sang đây; nó bắt đầu ở đây. Chỗ hổng của chương này vì thế không phải chỗ thiếu
bằng chứng mà là chỗ bằng chứng đã có không tới được hai khảo sát viết về cùng
đối tượng.
\end{itemize}
\end{tomtat}

\cauhoi
\begin{cauhoids}
  \item Mục~\ref{sec:ch17-loi-giai-thich-tu-viet} nêu bốn hệ quả của việc vế trái
đổi chủ, và cả bốn đều được viết ra như những chỗ làm việc kiểm một
\tn{lời giải thích tự sinh}{self-explanation}
khó hơn kiểm một bản đồ nhiệt. Hãy chọn một trong bốn và lập luận theo chiều
ngược lại: nêu một điều mà người kiểm \emph{làm được} với một chuỗi suy luận và
không làm được với một bản đồ nhiệt, đúng vì hệ quả ấy. Nếu bạn cho rằng hệ quả
mình chọn không cho lợi thế nào thì nói vì sao.
  \item Mục~\ref{sec:ch17-bai-mo-duong} đọc phần trong ngoặc của tính chất thứ
hai như một lời rào về độ trung thực. Hãy viết lại tính chất thứ hai ấy \emph{bỏ}
phần trong ngoặc đi, rồi cho biết câu còn lại tuyên bố điều gì mà câu đầy đủ
không tuyên bố, và mục~\ref{sec:ch17-nganh-doc-nguoc-lai} tìm thấy câu rút gọn ấy
ở đâu.
  \item Bài mở đường dùng chữ \enquote{faithfulness} đúng hai lần trong 43 trang
và một lần là trong dãy các năng lực trồi lên theo quy mô. Hãy đọc câu ấy trong
ngữ cảnh của nó và cho biết nghĩa ở đó gần với dòng nào của
bảng~\ref{tab:ch01-bon-quan-he} nhất, hoặc vì sao nó không gần dòng nào.
  \item Mục~\ref{sec:ch17-cho-dat-do-trung-thuc} xếp ba điều kiện của khảo sát
vào ba chỗ khác nhau và nói điều kiện thứ hai không rơi vào dòng nào của
bảng~\ref{tab:ch01-bon-quan-he}. Hãy dựng một ví dụ, có thể trên hàm chạy suốt
sách, trong đó một lời giải thích thỏa điều kiện thứ hai mà vi phạm điều kiện thứ
nhất, rồi nói ví dụ ấy chứng minh điều gì về việc ghép ba điều kiện bằng chữ
\enquote{và}.
  \item Bảng~\ref{tab:ch17-saliency} có mười hai giá trị. Hãy cho biết cần đổi
đúng một giá trị nào, thành số nào, để thứ hạng theo cột ngữ cảnh dài trở nên đơn
điệu giảm theo số tham số, rồi cho biết việc phải đổi một giá trị để cứu một câu
kết luận nói lên điều gì về câu kết luận ấy.
  \item Mục~\ref{sec:ch17-bai-neo-tro-ve} dựng một chuỗi ba bước trên
$f(x_1,x_2) = x_1 + 2x_2 + x_1x_2$ tại $(1,1)$, và cho thấy chuỗi ấy có thể có
mọi bước trung thực mà cả chuỗi không trung thực. Hãy dựng chuỗi đi theo chiều
ngược lại trên cùng hàm ấy tại $(0,1)$: một chuỗi \emph{không} thỏa điều kiện thứ
hai của định nghĩa~\ref{def:ch09-chuoi-trung-thuc} nhưng đường suy luận trong nó
lại là đường mô hình đã đi. Nếu bạn cho rằng không dựng được thì hãy nói vì sao.
  \item Hình 27 của khảo sát 2026 định nghĩa độ trung thực bằng câu hỏi dấu vết
có gây ra đáp án hay không, và liệt kê ba họ dụng cụ dưới câu ấy. Hãy cho biết,
với mỗi họ trong ba họ, nó đo được tính chất nào trong hai tính chất mà
mục~\ref{sec:ch09-dinh-nghia-lai} tách ra, rồi cho biết có họ nào trong ba đo
được điều kiện thứ nhất của định nghĩa~\ref{def:ch09-chuoi-trung-thuc} hay không.
  \item Mục~\ref{sec:ch17-chieu-cua-phep-chuyen} đặt điều kiện thứ năm lên một
dụng cụ đo tương lai và nói nó khác bốn điều kiện trước ở chỗ không đòi công bố
thêm gì. Hãy chọn một trong tám chỉ số mà mục~\ref{sec:ch09-ket-qua} liệt kê và
cho biết dưới điều kiện thứ năm thì chỉ số ấy phải khai điều gì, rồi cho biết
việc khai điều ấy có đổi điểm AUROC của nó trong bảng~\ref{tab:ch09-auroc} hay
không.
  \item Đọc mục 2.1 và mục 6, phần \enquote{Monitorability}, trong
\enquote{Faithfulness Metrics Don't Measure Faithfulness: A Meta-Evaluation with
Ground Truth} của Gur-Arieh và cộng sự~\autocite{p21metrics}. Hãy phát biểu đề
nghị thay độ trung thực bằng \tn{khả năng giám sát}{monitorability}
dưới dạng một câu nêu rõ nó bỏ
điều gì và giữ điều gì, rồi cho biết bảy chương của Phần~IV có làm cho đề nghị ấy
mạnh lên hay yếu đi, và vì sao.
\end{cauhoids}
